% Canonical LaTeX chapter; migrated from docs/design/engagement-audit.md.
\hypertarget{ascend-making-the-next-challenge-worth-opening}{%
\subsection{Ascend: making the next challenge worth
opening}\label{ascend-making-the-next-challenge-worth-opening}}

Audit and implementation: 2026-11-01. This revises the original UI brief
in response to the owner\textquotesingle s experience that the
implemented interface and tracks were not motivating. The four
destinations and shared CLI/server architecture remain useful. The
earlier restriction to a logging-first hero and extremely restrained
feedback is superseded by the challenge-first flow described here.

\hypertarget{what-the-reference-platforms-actually-contribute}{%
\subsubsection{What the reference platforms actually
contribute}\label{what-the-reference-platforms-actually-contribute}}

The audit used public first-party pages, not authenticated accounts or a
usability study. These observations describe published product flows;
they do not establish that the same changes will improve this
user\textquotesingle s long-term motivation.

\begin{longtable}[]{@{}P{.24\textwidth}P{.35\textwidth}P{.33\textwidth}@{}}
\toprule\noalign{}
Reference & Observed pattern & Ascend adaptation \\
\midrule\noalign{}
\endhead
\bottomrule\noalign{}
\endlastfoot
\href{https://blog.duolingo.com/new-duolingo-home-screen-design/}{Duolingo\textquotesingle s
learning-path design} & A guided next step, smaller units, distinctive
end-of-unit challenges, optional practice, and a way to return to the
current position. The article also explains practice interleaved into
the path. & Clearly marked current level, compact overview, numbered
connected challenges, distinct milestone nodes, jump-to-current
navigation. Optional practice stays separate; no spaced-repetition
algorithm is claimed. \\
\href{https://academy.hackthebox.com/}{HTB Academy} & Skill/job paths
built from modules; practical exercises, real-world examples, and
assessments. Difficulty and course scope are visible before starting. &
A challenge brief, estimated time, concrete finish line, direct resource
access, and a preview of the next milestone. Ascend does not provide
hosted lab machines or certifications. \\
\href{https://tryhackme.com/}{TryHackMe} & Beginner-friendly pathways,
guides and challenges, real-world scenarios, and exercises in lessons. &
Small investigations with a first action, optional progressive hints,
and an easy transition from doing the work to recording evidence. \\
\href{https://www.theodinproject.com/paths}{The Odin Project paths} & A
common foundation, ordered courses, and building real working projects.
& Retain deliberate prerequisite order while making the output of each
challenge concrete. \\
\end{longtable}

The public pages above were readable during the audit.
TryHackMe\textquotesingle s \texttt{/hacktivities} endpoint returned
429; an individual help article returned 403. No claim is made about
having inspected their signed-in dashboards. No platform assets,
proprietary course text, mascots, or course interfaces were copied.

\hypertarget{problems-in-the-previous-ascend-experience}{%
\subsubsection{Problems in the previous Ascend
experience}\label{problems-in-the-previous-ascend-experience}}

\begin{itemize}
\tightlist
\item
  Today offered \textbf{Log it} before showing a resource or helping the
  learner begin.
\item
  The completion control was hidden inside an optional-details
  disclosure.
\item
  Path showed a long, visually uniform list; the current level and next
  milestone were not prominent enough to create a sense of forward
  movement.
\item
  Passing structural pack validation was treated as a stronger
  content-quality signal than it is. The five imported cyber tracks had
  no explicit minute estimates. Several began with writing definitions,
  and one purported quick-win lane consisted of whole machines to root.
\item
  A file existing proves the presence of an artifact, not mastery. An
  installed tool\textquotesingle s version does not prove the
  corresponding exercise was completed.
\end{itemize}

\hypertarget{implemented-experience}{%
\subsubsection{Implemented experience}\label{implemented-experience}}

\textbf{Today:} one active-track overview, a clear current position,
honest completion counts, time choices, and one primary challenge. A
smaller optional challenge and milestone preview give an alternative and
a destination. The purple palette, consistent line icons, stronger type
hierarchy, connected route motif, and amber milestone treatment make
progress more legible. Dark mode and phone navigation use the same
system.

\textbf{Challenge:} Open challenge reveals the scenario, finish line,
optional concrete steps, resource link, and hints behind individual
disclosures. Existing packs without this metadata still show their
original completion criterion and resource. Hints never mark a challenge
done. There is no fake running lab or automatic timer.

\textbf{Save:} partial work and completion are explicit. Completion
calls the existing verification-backed logging endpoint. A failed check
saves the session and explains why the challenge is still open.
Confirmed completion shows a brief inline success panel and any newly
opened levels. It does not award invented XP or trigger an extra modal.
Completed challenges can be practiced again; manual mark-done is no
longer the prominent route around evidence checks.

\textbf{Path:} numbered level headers form a compact overview. The
current level opens by default; any other level can be expanded.
Connected circular challenge nodes have distinct ready, complete,
skipped, and locked states. Milestones use amber diamonds. Jump to
current level works with keyboard focus and reduced motion. Optional
drills are grouped below the main path.

\textbf{Stats:} earned milestone badges derive from actual \texttt{done}
nodes in the active track. Skipped nodes do not earn badges or inflate
the completed-step percentage. The existing streak, heatmap, minutes,
and note search remain available.

\hypertarget{content-changes-and-limits}{%
\subsubsection{Content changes and
limits}\label{content-changes-and-limits}}

The five cyber tracks are version 2, retain their existing node IDs, and
are now included in the bundled community registry. Their opening levels
use practical experiments and investigations with \texttt{brief},
\texttt{steps}, and \texttt{hints} metadata. Short optional discoveries
supply an actual small-session alternative. Larger exercises remain
larger; they are not relabeled as ten-minute tasks.

The first networking capstone can use a staged disposable-lab fault, so
completion no longer depends on waiting for a real outage. The nmap
opening checks the saved exercise evidence instead of whether nmap is
installed. A directory-as-file check, incorrect ROP challenge name, and
off-by-one drill descriptions were corrected. The later subject coverage
is still a draft curriculum, not a complete authored course or proof of
proficiency. File checks still check files, not understanding.

No XP economy, leaderboard, urgency prompts, external telemetry,
automatic trust, new accounts, or subscription gates are introduced. The
improvement to motivation must be judged through actual use. A useful
follow-up is whether the owner opens a challenge, reaches a concrete
result, and returns willingly over a week.
